% Un message empoisonné : remis indéfiniment, ou écarté après N livraisons vers
% des files de nouvelle tentative puis une file de lettres mortes.
\documentclass[border=6pt]{standalone}
\usepackage{sd}
\begin{document}
\begin{tikzpicture}[x=1mm, y=1mm, note/.append style={font=\sffamily\small},
    q/.style={draw=#1, fill=#1Bg, line width=0.7pt, rounded corners=2pt, minimum width=26mm, minimum height=8mm, font=\sffamily\footnotesize, align=center}]
  \node[pillo=Dred] at (45,64) {sans limite};
  \node[q=Dblue] (f1) at (10,50) {file};
  \node[soft=Dslate, minimum width=24mm, minimum height=8mm] (c1) at (70,50) {consommateur};
  \draw[arr=Dslate] (f1) -- node[lbl, above] {message empoisonné} (c1);
  \draw[arr=Dred] (c1.south) to[out=-150, in=-30] node[lblc=Dred, below] {échec, remis dans la file} (f1.south);
  \node[note, anchor=west, text=Dred, align=left] at (88,50) {remis sans fin ;\\les autres attendent derrière\\ou passent entre deux échecs};
  \begin{scope}[yshift=-44mm]
    \node[pillo=Dgreen] at (45,64) {avec limite et lettres mortes};
    \node[q=Dblue] (f) at (0,50) {file principale};
    \node[soft=Dslate, minimum width=22mm, minimum height=8mm] (c) at (38,50) {consommateur};
    \node[q=Dorange] (r1) at (78,58) {nouvelle tentative\\après 1 min};
    \node[q=Dorange] (r2) at (78,42) {nouvelle tentative\\après 10 min};
    \node[q=Dred] (dlq) at (118,50) {lettres mortes\\(DLQ)};
    \draw[arr=Dslate] (f) -- (c);
    \draw[arr=Dorange] (c) -- node[lblc=Dorange, above, sloped] {échec 1} (r1.west);
    \draw[arr=Dorange] (r1.south) -- node[lblc=Dorange, right] {échec 2} (r2.north);
    \draw[arr=Dred] (r2.east) -- node[lblc=Dred, below, sloped] {échec 3} (dlq.west);
    \node[note, anchor=north, align=center] at (58,32) {la file principale continue ; le message écarté est examiné,\\corrigé ou rejoué à la main depuis la DLQ};
  \end{scope}
\end{tikzpicture}
\end{document}
